%% toc: P3 39 39.5
\chapter{Games for Existential Second-Order Expressibility}
\label{sec:p03-w03-p03-058-eso-games}

First-order inexpressibility on finite structures is usually proved with the
Ehrenfeucht--Fra\"iss\'e game: two structures that the game cannot separate
in \(r\) rounds satisfy the same sentences of quantifier rank at most \(r\).
This chapter extends the method to existential second-order logic. The
extension adds a first stage in which relations of prescribed arities are
chosen on both structures; only then is the first-order game played on the
expanded structures. A second form of the game, the Ajtai--Fagin game, lets
Duplicator choose the second structure after Spoiler has chosen relations on
the first, and takes the query itself as a parameter. Both forms are sound
and complete for the corresponding non-definability claims. The chapter
proves these correspondences, works through a play of the relation-selection
stage, proves that connectivity of finite graphs is not definable in monadic
existential second-order logic by an explicit Duplicator strategy, and records
what the method has not achieved.

\emph{Inputs.} A vocabulary, a class \(\mathcal C\) of structures closed
under isomorphism, a Boolean query \(Q\) on \(\mathcal C\), and either a
sequence of arities with a number of rounds or a number of unary relations
with a number of rounds. \emph{Conclusions.} Either an existential
second-order sentence defining \(Q\) on \(\mathcal C\), or a proof that none
exists, obtained from winning strategies of Duplicator.

The presentation follows Section~5 of Kolaitis's manuscript \emph{On the
Expressive Power of Logics on Finite Models}, dated 1 August 2003
\cite[pp.~26--29]{Kolaitis2003ExpressivePowerNotes}. That manuscript is a
version of his chapter of the same title in the 2007 volume \emph{Finite
Model Theory and Its Applications}; the presentation follows the manuscript,
not the printed chapter, and every locator below refers to the manuscript and
its own numbering of sections and results. Kolaitis states the correspondence
theorems as methods and proves the connectivity theorem in outline. The
proofs in the annex are written out in full, and the places where they
depart from his outline are noted. Bold letters \(\mathbf A,\mathbf B\) for
structures and the names Spoiler and Duplicator are his.

\section{Vocabularies, definability and the first-order game}
\label{sec:p03-w03-p03-058-setting}

\begin{hbconvention}{def:p03-w03-p03-058-setting}{Vocabularies, existential second-order sentences and definability}
A \emph{vocabulary} \(\sigma\) is a finite set
\(\{R_1,\ldots,R_m,c_1,\ldots,c_s\}\) of relation symbols of positive
arities and constant symbols. A \(\sigma\)-structure \(\mathbf A\) has a
nonempty universe \(A\), finite or infinite
\cite[p.~3]{Kolaitis2003ExpressivePowerNotes}. A \emph{class} of
\(\sigma\)-structures is closed under isomorphism, and a \emph{Boolean
query} on a class \(\mathcal C\) is an isomorphism-invariant map
\(Q\colon\mathcal C\to\{0,1\}\) \cite[Def.~2.1, p.~4]{Kolaitis2003ExpressivePowerNotes}.
\begin{enumerate}
\item An \emph{ESO-sentence} has the form
      \((\exists P_1)\cdots(\exists P_k)\psi\) with \(k\ge0\), where the
      \(P_i\) are relation variables of positive arities and \(\psi\) is a
      first-order sentence over \(\sigma\cup\{P_1,\ldots,P_k\}\). It is
      \emph{monadic} when every \(P_i\) is unary
      \cite[p.~6]{Kolaitis2003ExpressivePowerNotes}. For relational
      \(\sigma\) these are the sentences of
      \cref{def:p03-existential-second-order}.
\item \(\mathbf A\entails(\exists P_1)\cdots(\exists P_k)\psi\) holds when
      there are relations \(S_1,\ldots,S_k\) on \(A\), \(S_i\) of the arity
      of \(P_i\), with
      \((\mathbf A,S_1,\ldots,S_k)\entails\psi\); here
      \((\mathbf A,S_1,\ldots,S_k)\) is the expansion of \(\mathbf A\) that
      interprets \(P_i\) by \(S_i\).
\item \(Q\) is \emph{ESO-definable} (\emph{monadic ESO-definable}) on
      \(\mathcal C\) if some ESO-sentence (monadic ESO-sentence) \(\Psi\)
      satisfies \(Q(\mathbf A)=1\iff\mathbf A\entails\Psi\) for every
      \(\mathbf A\in\mathcal C\) \cite[Def.~2.3, p.~5]{Kolaitis2003ExpressivePowerNotes}.
\item For \(r\ge0\), \(\mathbf A\equiv_r\mathbf B\) means that
      \(\mathbf A\) and \(\mathbf B\) satisfy the same first-order sentences
      of quantifier rank at most \(r\), quantifier rank being that of
      \cref{def:p03-quantifier-rank}.
\item In the \emph{\(r\)-move Ehrenfeucht--Fra\"iss\'e game} on two
      \(\sigma\)-structures, Spoiler and Duplicator alternate as in
      \cref{def:p03-ef-game}; if \((a_1,b_1),\ldots,(a_r,b_r)\) are the
      selected pairs, Duplicator wins the run when the relation
      \[
       \{(a_i,b_i):1\le i\le r\}\cup\{(c_j^{\mathbf A},c_j^{\mathbf B}):
       1\le j\le s\}
      \]
      is a partial isomorphism, that is, the graph of an isomorphism between
      the substructures generated by \(\{a_1,\ldots,a_r\}\) and by
      \(\{b_1,\ldots,b_r\}\) \cite[Def.~3.1, p.~9]{Kolaitis2003ExpressivePowerNotes}.
      Duplicator \emph{wins the game} when she has a winning strategy.
\end{enumerate}
A finite set \(p\subseteq A\times B\) is a partial isomorphism when
\(a=a'\iff b=b'\) for \((a,b),(a',b')\in p\) and every relation symbol
holds of a tuple of first coordinates exactly when it holds of the
corresponding tuple of second coordinates. With constants present, the
constant pairs are part of every position from the start.
\index{existential second-order logic!monadic}\index{Boolean query}
\end{hbconvention}

For finite relational signatures, the first-order game from a position
with parameters is characterized in \cref{thm:p03-w03-p03-038-pointed-ef}.
The following lemma carries that characterization over to vocabularies with
constants, which Kolaitis's games allow.

\begin{hblemma}{thm:p03-w03-p03-058-ef-constants}{proved=annex:p03-w03-p03-058-ef-constants}{The first-order game with constants}
Let \(\sigma\) be a vocabulary, \(r\ge0\), and \(\mathbf A,\mathbf B\)
\(\sigma\)-structures.
\begin{enumerate}
\item Duplicator wins the \(r\)-move game on \(\mathbf A\) and \(\mathbf B\)
      if and only if \(\mathbf A\equiv_r\mathbf B\).
\item There are finitely many \(\equiv_r\)-classes of \(\sigma\)-structures,
      and for each \(\sigma\)-structure \(\mathbf A\) there is a first-order
      \(\sigma\)-sentence \(\varphi^{\mathbf A}_r\) of quantifier rank at
      most \(r\) such that \(\mathbf B\entails\varphi^{\mathbf A}_r\) if and only
      if \(\mathbf A\equiv_r\mathbf B\). The set of all sentences
      \(\varphi^{\mathbf A}_r\) is finite.
\end{enumerate}
\end{hblemma}

This is Kolaitis's Theorem~3.8 together with his notation
\(\varphi^{\mathbf A}_r\) for the sentence describing the class of
\(\mathbf A\) \cite[Thm.~3.8, p.~11; Def.~3.13, p.~13]{Kolaitis2003ExpressivePowerNotes}.
Both items hold for infinite structures; the finiteness of \(\sigma\) is what
makes the set of describing sentences finite.

\section{The Ehrenfeucht--Fra\"iss\'e game for existential second-order logic}
\label{sec:p03-w03-p03-058-eso-game}

\begin{hbdefinition}{def:p03-w03-p03-058-eso-game}{The \((\langle s_1,\ldots,s_k\rangle,r)\)-Ehrenfeucht--Fra\"iss\'e game}
Let \(s_1,\ldots,s_k,r\) be positive integers, \(\sigma\) a vocabulary and
\(\mathbf A,\mathbf B\) two \(\sigma\)-structures. A run of the
\((\langle s_1,\ldots,s_k\rangle,r)\)-Ehrenfeucht--Fra\"iss\'e game on
\(\mathbf A\) and \(\mathbf B\) has two stages.
\begin{enumerate}
\item \emph{Relation selection.} Spoiler chooses relations
      \(S_1,\ldots,S_k\) on \(A\), \(S_i\) of arity \(s_i\). Duplicator then
      chooses relations \(S'_1,\ldots,S'_k\) on \(B\) of the same arities.
      Any relations of the right arities may be chosen, including empty
      ones.
\item \emph{First-order stage.} The players play a run of the \(r\)-move
      Ehrenfeucht--Fra\"iss\'e game on the expanded structures
      \((\mathbf A,S_1,\ldots,S_k)\) and \((\mathbf B,S'_1,\ldots,S'_k)\),
      whose vocabulary is \(\sigma\) together with new relation symbols
      \(P_1,\ldots,P_k\) of arities \(s_1,\ldots,s_k\).
\end{enumerate}
Duplicator wins the run when the selected pairs, together with the pairs of
constants, form a partial isomorphism between the expanded structures; so
the chosen relations, not only those of \(\sigma\), must be preserved and
reflected. Duplicator \emph{wins the game} when she has a winning strategy:
for every choice of \(S_1,\ldots,S_k\) on \(A\) there are
\(S'_1,\ldots,S'_k\) on \(B\) such that Duplicator wins the \(r\)-move
game on \((\mathbf A,S_1,\ldots,S_k)\) and \((\mathbf B,S'_1,\ldots,S'_k)\)
\cite[Def.~5.1, p.~26]{Kolaitis2003ExpressivePowerNotes}.
\index{Ehrenfeucht--Fraisse game@Ehrenfeucht--Fra\"iss\'e game!for existential second-order logic}
\end{hbdefinition}

Both structures are fixed before the run begins, and the relation-selection
stage happens once, before any element is chosen. Spoiler's choice is a
guess of witnesses on \(\mathbf A\); Duplicator must copy it on
\(\mathbf B\) well enough to survive \(r\) further rounds. With \(k=0\) the
game is the first-order game of \cref{def:p03-w03-p03-058-setting}. The game
is not symmetric: Spoiler selects relations only on \(\mathbf A\).

\begin{hbwarning}{def:p03-w03-p03-058-not-if}{Two different second-order games}
The relation-selection game compares two structures and involves no
sentence: its outcome says which sentences cannot separate the structures.
The semantic evaluation game of independence-friendly logic
(\cref{def:p03-if-logic}) is played on one structure and one sentence, and
its outcome says whether that sentence is true there. In the evaluation
game a player chooses elements, and a slash set restricts the information on
which a choice may depend; in the comparison game the second-order content
is an explicit choice of whole relations. Independence-friendly and
dependence-logic sentences have exactly the expressive power of
existential second-order sentences (\cref{thm:p03-dependence-eso}), so a
query shown here not to be ESO-definable is also not definable by such a
sentence; the two games are nevertheless different games and answer
different questions.
\end{hbwarning}

\begin{hbtheorem}{thm:p03-w03-p03-058-preservation}{proved=annex:p03-w03-p03-058-preservation}{Duplicator wins preserve existential second-order sentences}
Let \(\Psi\) be an ESO-sentence \((\exists P_1)\cdots(\exists P_k)\psi\),
where \(P_i\) has arity \(s_i\) and \(\psi\) is a first-order sentence of
quantifier rank at most \(r\), with \(k,r\ge1\). If
\(\mathbf A\entails\Psi\) and Duplicator wins the
\((\langle s_1,\ldots,s_k\rangle,r)\)-Ehrenfeucht--Fra\"iss\'e game on
\(\mathbf A\) and \(\mathbf B\), then \(\mathbf B\entails\Psi\).
\end{hbtheorem}

This is Kolaitis's Proposition~5.2, stated there for a matrix of quantifier
rank \(r\) \cite[p.~26]{Kolaitis2003ExpressivePowerNotes}. The hypotheses
\(k,r\ge1\) cost nothing: a vacuous quantifier \((\exists P)\) can be added
in front of any ESO-sentence, and a sentence of rank at most \(r\) also
has rank at most \(r+1\).

\begin{hbtheorem}{thm:p03-w03-p03-058-eso-method}{proved=annex:p03-w03-p03-058-eso-method}{Soundness and completeness of the game method for ESO}
Let \(\sigma\) be a vocabulary, \(\mathcal C\) a class of
\(\sigma\)-structures and \(Q\) a Boolean query on \(\mathcal C\).
\begin{enumerate}
\item \(Q\) is not ESO-definable on \(\mathcal C\) if and only if, for every
      sequence of positive integers \(s_1,\ldots,s_k,r\), there are
      \(\mathbf A,\mathbf B\in\mathcal C\) with \(Q(\mathbf A)=1\),
      \(Q(\mathbf B)=0\), such that Duplicator wins the
      \((\langle s_1,\ldots,s_k\rangle,r)\)-Ehrenfeucht--Fra\"iss\'e game on
      \(\mathbf A\) and \(\mathbf B\).
\item \(Q\) is not monadic ESO-definable on \(\mathcal C\) if and only if,
      for all positive integers \(k,r\), there are such \(\mathbf A\) and
      \(\mathbf B\) on which Duplicator wins the
      \((\langle1,\ldots,1\rangle,r)\)-game with \(k\) entries \(1\).
\end{enumerate}
\end{hbtheorem}

The ``if'' directions are the soundness half of Kolaitis's Method~5.3 and
the ``only if'' directions its completeness half, which he states without
proof \cite[Method~5.3, pp.~26--27]{Kolaitis2003ExpressivePowerNotes}. The
second item is the monadic restriction he singles out as the case in which
the method has succeeded \cite[p.~27]{Kolaitis2003ExpressivePowerNotes}.
The class \(\mathcal C\) is arbitrary: it may contain infinite structures.

\begin{hbexample}{def:p03-w03-p03-058-c4-c5}{What the relation-selection stage adds}
Let \(C_n\) be the undirected cycle with vertices \(0,\ldots,n-1\) and
edges \(\{i,i+1\}\) (indices modulo \(n\)), so that \(E\) is symmetric and
irreflexive.

\emph{First-order stage alone.} Duplicator wins the \(2\)-move game on
\(C_4\) and \(C_5\). After any first pair \((a_1,b_1)\), whatever Spoiler
plays next stands to the element already selected in that structure in one
of three ways: equal, adjacent, or distinct and nonadjacent. In both cycles
every vertex has an equal partner (itself), an adjacent one and a distinct
nonadjacent one, so Duplicator answers with an element standing to the
element selected in the other structure in the same way. Both selected
pairs then agree on equality and on \(E\) in both directions, which is all
a partial isomorphism requires here.

\emph{With one unary relation.} Spoiler wins the
\((\langle1\rangle,2)\)-game on \(C_4\) and \(C_5\). He chooses
\(S=\{0,2\}\) on \(C_4\); then every edge of \(C_4\) joins a member of
\(S\) to a nonmember. Whatever \(S'\subseteq\{0,\ldots,4\}\) Duplicator
chooses, some edge \(\{u,u+1\}\) of \(C_5\) has both ends in \(S'\) or both
outside it: otherwise membership in \(S'\) would alternate around the cycle,
and an alternating labelling of a cycle of odd length returns to its start
with the opposite label. Spoiler plays \(u\) and then \(u+1\) in \(C_5\).
Duplicator's answers \(a_1,a_2\) in \(C_4\) must be adjacent, to reflect
\(E(u,u+1)\), and then exactly one of them lies in \(S\), whereas \(u\)
and \(u+1\) agree on \(S'\). The selected pairs are not a partial
isomorphism of the expanded structures. This matches
\cref{thm:p03-w03-p03-058-preservation}: the monadic sentence
\((\exists R)(\forall x)(\forall y)\bigl(E(x,y)\to(R(x)\leftrightarrow\neg
R(y))\bigr)\), Kolaitis's definition of 2-colorability
\cite[Example~2.5, p.~7]{Kolaitis2003ExpressivePowerNotes}, has a matrix of
quantifier rank \(2\), holds in \(C_4\) and fails in \(C_5\).
\end{hbexample}

The game method also restates a central open problem of complexity theory.
For the next result, \(\mathcal G\) is the class of finite graphs, that is,
of finite structures \((V,E)\) with one binary relation, and for a Boolean
query \(Q\) on \(\mathcal G\) the set \(E(Q)\) is the set of encodings of
the graphs with \(Q=1\) under
\cref{def:p02-w02-p02-036-fagin-encoding}. The complexity classes and
reductions are those of \cref{def:p02-complexity-standing}, a word over
\(\{0,1,\#\}\) being read as a binary word through the code
\(0\mapsto00\), \(1\mapsto01\), \(\#\mapsto10\).

\begin{hbcorollary}{thm:p03-w03-p03-058-np-conp}{proved=annex:p03-w03-p03-058-np-conp}{\(\mathsf{NP}\ne\mathsf{coNP}\) as a statement about games}
Let \(Q\) be a Boolean query on \(\mathcal G\) such that \(E(Q)\) is
\(\mathsf{NP}\)-complete. Then the following are equivalent.
\begin{enumerate}
\item \(\mathsf{NP}\ne\mathsf{coNP}\).
\item For every sequence of positive integers \(s_1,\ldots,s_k,r\) there are
      finite graphs \(\mathbf G\) and \(\mathbf H\) with
      \(Q(\mathbf G)=0\) and \(Q(\mathbf H)=1\) such that Duplicator wins
      the \((\langle s_1,\ldots,s_k\rangle,r)\)-Ehrenfeucht--Fra\"iss\'e
      game on \(\mathbf G\) and \(\mathbf H\).
\end{enumerate}
\end{hbcorollary}

Kolaitis states this equivalence for 3-colorability, deriving it from
Fagin's theorem and the game method
\cite[Thm.~4.7 and Cor.~4.8, p.~25; Cor.~5.4, p.~27]{Kolaitis2003ExpressivePowerNotes}.
The corollary above is that statement with the property of
3-colorability that it uses, \(\mathsf{NP}\)-completeness, made a
hypothesis; the next theorem supplies that property, and the instance follows.
Fagin's theorem enters as \cref{thm:p03-fagin}; its proof in
Fagin's original form is \cref{thm:p09-fagin74-theorem6}.

\begin{hbdefinition}{def:p03-w03-p03-058-three-col}{3-colorability}
A graph \(\mathbf G=(V,E)\), with \(E\) an arbitrary binary relation, is
\emph{3-colorable} when some map \(c\colon V\to\{1,2,3\}\) satisfies
\(c(u)\ne c(v)\) whenever \(E(u,v)\). A graph with a loop \(E(u,u)\) is
therefore not 3-colorable. \(\mathrm{3COL}\) is the Boolean query on
\(\mathcal G\) with \(\mathrm{3COL}(\mathbf G)=1\) exactly when
\(\mathbf G\) is 3-colorable. Kolaitis defines it on all finite graphs by
a monadic ESO-sentence with two quantified unary variables, the third color
being the complement of the other two
\cite[Example~2.5, p.~7]{Kolaitis2003ExpressivePowerNotes}.
\end{hbdefinition}

\begin{hbtheorem}{thm:p03-w03-p03-058-three-col-np}{proved=annex:p03-w03-p03-058-three-col-np}{3-colorability is \(\mathsf{NP}\)-complete}
\(E(\mathrm{3COL})\) is \(\mathsf{NP}\)-complete under polynomial-time
many-one reductions, words over \(\{0,1,\#\}\) being read as binary words
through the code \(0\mapsto00\), \(1\mapsto01\), \(\#\mapsto10\).
\end{hbtheorem}

The proof reduces satisfiability (\cref{thm:p02-sat-np-complete}) to
3-colorability through the definitional clause form of
\cref{thm:p02-tseitin}, whose clauses have at most three literals. Each
clause is replaced by two copies of a six-edge gadget, which can be colored
with its output in the color class of ``true'' only when one of its inputs is
there.

\begin{hbcorollary}{thm:p03-w03-p03-058-np-conp-3col}{proved=annex:p03-w03-p03-058-np-conp-3col}{Kolaitis's game form of \(\mathsf{NP}\ne\mathsf{coNP}\)}
\(\mathsf{NP}\ne\mathsf{coNP}\) if and only if, for every sequence of
positive integers \(s_1,\ldots,s_k,r\), there are finite graphs
\(\mathbf G\) and \(\mathbf H\) such that \(\mathbf G\) is not 3-colorable,
\(\mathbf H\) is 3-colorable, and Duplicator wins the
\((\langle s_1,\ldots,s_k\rangle,r)\)-Ehrenfeucht--Fra\"iss\'e game on
\(\mathbf G\) and \(\mathbf H\).
\end{hbcorollary}

This is Kolaitis's Corollary~5.4
\cite[p.~27]{Kolaitis2003ExpressivePowerNotes}.

\section{Ajtai--Fagin games}
\label{sec:p03-w03-p03-058-af}

In the game of \cref{def:p03-w03-p03-058-eso-game} Duplicator in effect
commits to \(\mathbf B\) before Spoiler reveals his relations on
\(\mathbf A\). Kolaitis observes that the choice of the two structures can
itself be made part of the play, and that Ajtai and Fagin introduced a
variant in which Duplicator chooses \(\mathbf B\) only after seeing
Spoiler's relations \cite[p.~27]{Kolaitis2003ExpressivePowerNotes}.

\begin{hbdefinition}{def:p03-w03-p03-058-af-game}{Ajtai--Fagin games}
Let \(\mathcal C\) be a class of \(\sigma\)-structures, \(Q\) a Boolean query
on \(\mathcal C\), and \(s_1,\ldots,s_k,r\) positive integers. A run of the
\((\langle s_1,\ldots,s_k\rangle,r)\)-\emph{Ajtai--Fagin game for \(Q\) on
\(\mathcal C\)} proceeds as follows.
\begin{enumerate}
\item Duplicator chooses \(\mathbf A\in\mathcal C\) with \(Q(\mathbf A)=1\).
\item Spoiler chooses relations \(S_1,\ldots,S_k\) on \(A\), \(S_i\) of
      arity \(s_i\).
\item Duplicator chooses \(\mathbf B\in\mathcal C\) with
      \(Q(\mathbf B)=0\) and then relations \(S'_1,\ldots,S'_k\) on \(B\),
      \(S'_i\) of arity \(s_i\).
\item The players play the \(r\)-move Ehrenfeucht--Fra\"iss\'e game on
      \((\mathbf A,S_1,\ldots,S_k)\) and \((\mathbf B,S'_1,\ldots,S'_k)\),
      with the winning condition of \cref{def:p03-w03-p03-058-eso-game}.
\end{enumerate}
A player who cannot move as required loses. Duplicator \emph{wins the game}
when there is \(\mathbf A\in\mathcal C\) with \(Q(\mathbf A)=1\) such that,
for all \(S_1,\ldots,S_k\) on \(A\) of arities \(s_1,\ldots,s_k\), there
are \(\mathbf B\in\mathcal C\) with \(Q(\mathbf B)=0\) and
\(S'_1,\ldots,S'_k\) on \(B\) such that Duplicator wins the \(r\)-move game
on \((\mathbf A,S_1,\ldots,S_k)\) and \((\mathbf B,S'_1,\ldots,S'_k)\).
When all \(s_i\) equal \(1\) this is the \emph{\((k,r)\)-Ajtai--Fagin game
for \(Q\) on \(\mathcal C\)}, Kolaitis's Definition~5.5
\cite[pp.~27--28]{Kolaitis2003ExpressivePowerNotes}.
\index{Ajtai--Fagin game}
\end{hbdefinition}

Two features distinguish this game from
\cref{def:p03-w03-p03-058-eso-game}. The structures are moves, chosen inside
\(\mathcal C\) with prescribed values of \(Q\), so the query is a parameter of
the game rather than something the game is used to test; and
\(\mathbf B\) may depend on Spoiler's relations. Kolaitis gives the monadic
form and remarks that the extension to arbitrary arities needs only changes
of notation \cite[p.~27]{Kolaitis2003ExpressivePowerNotes}; the definition
above makes that extension explicit.

\begin{hbtheorem}{thm:p03-w03-p03-058-af-method}{proved=annex:p03-w03-p03-058-af-method}{Soundness and completeness of Ajtai--Fagin games}
Let \(\mathcal C\) be a class of \(\sigma\)-structures and \(Q\) a Boolean
query on \(\mathcal C\).
\begin{enumerate}
\item \(Q\) is not monadic ESO-definable on \(\mathcal C\) if and only if,
      for all positive integers \(k\) and \(r\), Duplicator wins the
      \((k,r)\)-Ajtai--Fagin game for \(Q\) on \(\mathcal C\).
\item \(Q\) is not ESO-definable on \(\mathcal C\) if and only if, for every
      sequence of positive integers \(s_1,\ldots,s_k,r\), Duplicator wins
      the \((\langle s_1,\ldots,s_k\rangle,r)\)-Ajtai--Fagin game for \(Q\)
      on \(\mathcal C\).
\item For fixed \(s_1,\ldots,s_k,r\): if there are
      \(\mathbf A,\mathbf B\in\mathcal C\) with \(Q(\mathbf A)=1\),
      \(Q(\mathbf B)=0\) on which Duplicator wins the
      \((\langle s_1,\ldots,s_k\rangle,r)\)-Ehrenfeucht--Fra\"iss\'e game,
      then Duplicator wins the
      \((\langle s_1,\ldots,s_k\rangle,r)\)-Ajtai--Fagin game for \(Q\) on
      \(\mathcal C\).
\end{enumerate}
\end{hbtheorem}

The first item is Kolaitis's Method~5.6, soundness and completeness, which
he states without proof \cite[p.~28]{Kolaitis2003ExpressivePowerNotes}. The
third item makes precise in what sense the Ajtai--Fagin game is easier for
Duplicator: every strategy in the earlier game yields one in the new game.
By the first two items and \cref{thm:p03-w03-p03-058-eso-method}, Duplicator
wins all the Ajtai--Fagin games for \(Q\) exactly when, for all parameters,
she wins the earlier game on some pair of structures with \(Q\)-values
\(1\) and \(0\). For fixed parameters only the implication of the third
item is asserted.

\section{Connectivity is not monadic existential second-order}
\label{sec:p03-w03-p03-058-connectivity}

Kolaitis reports that Fagin first proved this with the games of
\cref{def:p03-w03-p03-058-eso-game}, and that Fagin, Stockmeyer and Vardi
later gave a much simpler proof with Ajtai--Fagin games and a counting
criterion for the first-order game
\cite[p.~28]{Kolaitis2003ExpressivePowerNotes}. The proof below follows the
second route. The counting criterion is needed only for graphs with
colours.

\begin{hbdefinition}{def:p03-w03-p03-058-neighbourhoods}{Coloured graphs, neighbourhoods and \(d\)-equivalence}
A \emph{\(k\)-coloured graph} is a structure
\(\mathbf A=(A,E,S_1,\ldots,S_k)\) with \(E\) symmetric and irreflexive and
each \(S_i\) unary. Its \emph{colour} at \(a\) is the vector
\((c_1,\ldots,c_k)\in\{0,1\}^k\) with \(c_i=1\) exactly when \(a\in S_i\).
The distance \(\operatorname{dist}(a,b)\) is the length of a shortest
\(E\)-path from \(a\) to \(b\), and \(\infty\) if there is none. For an
integer \(d\ge1\):
\begin{enumerate}
\item the \emph{neighbourhood of radius \(d\)} of \(a\) is
      \(N(a,d)=\{b\in A:\operatorname{dist}(a,b)<d\}\), so
      \(N(a,1)=\{a\}\);
\item the \emph{\(d\)-type} of \(a\) is the isomorphism type of the
      substructure of \((\mathbf A,a)\) with universe \(N(a,d)\), where
      \((\mathbf A,a)\) carries \(a\) as a distinguished element;
\item \(\mathbf A\) and \(\mathbf B\) are \emph{\(d\)-equivalent} when, for
      every \(d\)-type, they have the same number of elements of that
      \(d\)-type.
\end{enumerate}
For such structures \(E\) is the Gaifman graph, so these are Kolaitis's
neighbourhoods, \(d\)-types and \(d\)-equivalence specialized to coloured
graphs \cite[Defs.~3.22, 3.24 and~3.25, pp.~18--19]{Kolaitis2003ExpressivePowerNotes}.
For a tuple \(\bar a=(a_1,\ldots,a_j)\),
\(N(\bar a,d)=\bigcup_iN(a_i,d)\), and
\(\operatorname{dist}(x,\bar a)=\min_i\operatorname{dist}(x,a_i)\).
\index{neighbourhood!in a graph}\index{d-equivalence@\(d\)-equivalence}
\end{hbdefinition}

\begin{hblemma}{thm:p03-w03-p03-058-hanf}{proved=annex:p03-w03-p03-058-hanf}{Counting neighbourhood types wins the first-order game}
Let \(\mathbf A\) and \(\mathbf B\) be finite \(k\)-coloured graphs, \(r\ge1\)
and \(d\ge3^{r-1}\). If \(\mathbf A\) and \(\mathbf B\) are
\(d\)-equivalent, then Duplicator wins the \(r\)-move
Ehrenfeucht--Fra\"iss\'e game on \(\mathbf A\) and \(\mathbf B\).
\end{hblemma}

Kolaitis states this for arbitrary vocabularies as his Theorem~3.26,
attributed to Fagin, Stockmeyer and Vardi, with a hint that maintains
isomorphisms between unions of neighbourhoods of radius \(3^{r-j}\) after
\(j\) moves \cite[Thm.~3.26, p.~19]{Kolaitis2003ExpressivePowerNotes}. The
annex proof carries out that plan for finite coloured graphs, the case used
below; the restriction to coloured graphs makes every edge between two
points of a neighbourhood visible in the induced substructure. The general
statement is not proved here.

For the strategy, fix \(k,r\ge1\) and put
\[
 d=3^{r-1},\qquad M=2^{k(2d+1)},\qquad n=(M+1)(2d+1).
\]
For a colouring \(S_1,\ldots,S_k\) of the cycle \(C_n\) and a vertex \(i\),
the \emph{window} of \(i\) is the sequence of colours of
\(i-d,i-d+1,\ldots,i+d\) (indices modulo \(n\)); there are \(M\) possible
windows.

\begin{hblemma}{thm:p03-w03-p03-058-pinch}{proved=annex:p03-w03-p03-058-pinch}{Pinching a coloured cycle}
Let \(C_n\) carry colours \(S_1,\ldots,S_k\), and let
\(0\le p<q\le n-1\) be vertices with the same window such that
\(q-p\ge2d+1\) and \(n-(q-p)\ge2d+1\). Let \(\mathbf B\) have the vertices
and colours of \(C_n\) and the edges \(\{i,i+1\}\) for
\(i\notin\{p,q\}\) together with \(\{p,q+1\}\) and \(\{p+1,q\}\). Then
\begin{enumerate}
\item \(\mathbf B\) is the disjoint union of two cycles, on
      \(\{p+1,\ldots,q\}\) and on the remaining vertices, each of length at
      least \(2d+1\); in particular \(\mathbf B\) is not connected;
\item every vertex has the same \(d\)-type in \((C_n,S_1,\ldots,S_k)\) and
      in \((\mathbf B,S_1,\ldots,S_k)\), so the two coloured graphs are
      \(d\)-equivalent.
\end{enumerate}
\end{hblemma}

\begin{hbtheorem}{thm:p03-w03-p03-058-connectivity}{proved=annex:p03-w03-p03-058-connectivity}{Connectivity is not monadic ESO-definable}
A graph \((V,E)\) is \emph{connected} when any two vertices are joined by a
path in the symmetric closure of \(E\), and \emph{strongly connected} when
for any two vertices \(u,v\) there is a directed \(E\)-path from \(u\) to
\(v\). For all positive integers \(k\) and \(r\), Duplicator wins the
\((k,r)\)-Ajtai--Fagin game for connectivity and the one for strong
connectivity on the class of all finite graphs, and the game for
connectivity on the class of finite graphs with symmetric irreflexive edge
relation. Hence none of these three queries is monadic ESO-definable on its
class.
\end{hbtheorem}

This is Kolaitis's Theorem~5.7 \cite[pp.~28--29]{Kolaitis2003ExpressivePowerNotes}.
His sketch asks for two far-apart vertices of the same \(d\)-type and pinches
the cycle there. The annex proof asks instead for equal windows, which fix
the colours on each side of the two vertices separately: a \(d\)-type of a
vertex on an undirected cycle may be realized by a reflected colour pattern,
and pinching at two vertices whose patterns are reflections of each other
does not in general preserve \(d\)-types. The counting that supplies the
two vertices is the pigeonhole argument above.

\begin{hbexample}{def:p03-w03-p03-058-strategy}{Duplicator's strategy, move by move}
For given \(k\) and \(r\), Duplicator plays as follows in the
\((k,r)\)-Ajtai--Fagin game for connectivity.
\begin{enumerate}
\item She chooses \(\mathbf A=C_n\) with \(d,M,n\) as above. It is connected.
\item Spoiler colours \(C_n\) by \(S_1,\ldots,S_k\).
\item Among the \(M+1\) vertices \(0,2d+1,2(2d+1),\ldots,M(2d+1)\) two have
      the same window; call them \(p<q\). Then \(q-p\ge2d+1\), and
      \(n-(q-p)\ge n-M(2d+1)=2d+1\).
\item She chooses the pinched graph \(\mathbf B\) of
      \cref{thm:p03-w03-p03-058-pinch} and the same colours
      \(S'_i=S_i\). It is disconnected.
\item In the \(r\)-move game she follows the strategy of
      \cref{thm:p03-w03-p03-058-hanf}, which applies because the coloured
      graphs are \(d\)-equivalent and \(d=3^{r-1}\).
\end{enumerate}
For \(k=1\) and \(r=1\) the numbers are \(d=1\), \(M=8\) and \(n=27\); the
window of a vertex is the colours of the vertex and its two neighbours, and
in the one-move game Duplicator answers each vertex by the same vertex of
the other graph, which has the same colour because the two graphs share
vertices and colours.
\end{hbexample}

\begin{hbcorollary}{thm:p03-w03-p03-058-complement}{proved=annex:p03-w03-p03-058-complement}{Monadic ESO is not closed under complement}
Let \(\delta\) be the monadic ESO-sentence
\[
 (\exists S)\bigl((\exists x)S(x)\wedge(\exists y)\neg S(y)\wedge
 (\forall z)(\forall w)(S(z)\wedge\neg S(w)\to\neg E(z,w))\bigr),
\]
and let \(\delta'\) be \(\delta\) with \(\neg E(z,w)\) replaced by
\(\neg E(z,w)\wedge\neg E(w,z)\).
\begin{enumerate}
\item On finite graphs with symmetric irreflexive edge relation, \(\delta\)
      defines disconnectivity.
\item On all finite graphs, \(\delta\) defines the failure of strong
      connectivity and \(\delta'\) defines disconnectivity.
\end{enumerate}
In each case the complementary query is not monadic ESO-definable on the
class. Hence the monadic ESO-definable queries on finite graphs are not
closed under complement.
\end{hbcorollary}

The sentence \(\delta\) is Kolaitis's definition of disconnectivity
\cite[Example~2.5, p.~7]{Kolaitis2003ExpressivePowerNotes}, and the
corollary is his Corollary~5.8 \cite[p.~29]{Kolaitis2003ExpressivePowerNotes}.
His graphs are arbitrary binary relations \cite[p.~4]{Kolaitis2003ExpressivePowerNotes}.
On a graph that is not symmetric, \(\delta\) only forbids edges leaving
\(S\): the graph with vertices \(a,b\) and the single edge \((b,a)\) is
connected and satisfies \(\delta\) with \(S=\{a\}\). The second item
records what \(\delta\) defines there and the variant that defines
disconnectivity.
The complement of a monadic ESO query is defined by the negated sentence,
which is universal monadic second-order; so the corollary separates monadic
ESO from its universal dual on finite graphs.

\section{Reach and limits of the method}
\label{sec:p03-w03-p03-058-limits}

\begin{hbwarning}{def:p03-w03-p03-058-limits}{What the games have not achieved}
\Cref{thm:p03-w03-p03-058-eso-method,thm:p03-w03-p03-058-af-method} are
exact characterizations, and \cref{thm:p03-w03-p03-058-np-conp} turns
\(\mathsf{NP}\ne\mathsf{coNP}\) into a statement about the games. The
characterizations do not make the games easy to win. Kolaitis records that
the method has had limited success, because the combinatorics becomes
formidable as soon as some \(s_i\) exceeds \(1\), and that its successes
have been essentially confined to monadic ESO
\cite[pp.~27 and~29]{Kolaitis2003ExpressivePowerNotes}. The connectivity
theorem has no consequence for \(\mathsf{NP}\) versus \(\mathsf{coNP}\):
Kolaitis notes that connectivity is a polynomial-time computable query and
that monadic ESO does not define every \(\mathsf{NP}\) query, so a query
outside monadic ESO need not lie outside \(\mathsf{NP}\)
\cite[p.~29]{Kolaitis2003ExpressivePowerNotes}. In the 2003 manuscript,
Kolaitis poses as open the problem of finding an \(\mathsf{NP}\) query on
finite graphs that is not definable by any ESO-sentence all of whose
quantified relation variables are binary
\cite[Problem~5.9, p.~29]{Kolaitis2003ExpressivePowerNotes}. This is the
status reported in that source, not a claim about the problem's later status.
\end{hbwarning}

Right after posing that problem, Kolaitis records that the games have
nonetheless bounded monadic ESO on finite graphs that carry a
\emph{built-in} relation, namely a successor relation (a result he credits
to de Rougemont) and a total order (credited to Schwentick), and that such
results are regarded as the first step towards an analysis of definability
in binary existential second-order logic
\cite[p.~29]{Kolaitis2003ExpressivePowerNotes}. The two results are the
following theorems.

\begin{hbconvention}{def:p03-w03-p03-058-builtin}{Graphs with a built-in successor or order}
A \emph{finite structure with successor} is a finite structure
\(\mathbf G=(V,E,\mathrm{Suc})\) in which \(E\) is a binary relation on
\(V\) and \(\mathrm{Suc}\) is the graph of a successor function on \(V\).
A \emph{finite ordered graph} is a finite structure
\(\mathbf G=(V,E,\le)\) in which \(E\) is a binary relation on \(V\) and
\(\le\) is a linear order on \(V\)
\cite[p.~29 and p.~4]{Kolaitis2003ExpressivePowerNotes}. The
\emph{Connectivity query} takes the value \(1\) on \(\mathbf G\) when
\(\mathbf G\) is connected and \(0\) otherwise
\cite[Example~2.2, pp.~4--5]{Kolaitis2003ExpressivePowerNotes}; on these
structures connectedness concerns the edge relation \(E\), while
\(\mathrm{Suc}\) and \(\le\) are additional relations of every structure in
the class, which a defining sentence may mention and which neither player
chooses.
\end{hbconvention}

\begin{hbtheorem}{thm:p03-w03-p03-058-successor}{cited=deRougemont1987}{Connectivity with a built-in successor}
The Connectivity query is not monadic ESO-definable on the class of finite
structures with successor, that is, of finite structures
\(\mathbf G=(V,E,\mathrm{Suc})\) in which \(E\) is a binary relation on
\(V\) and \(\mathrm{Suc}\) is the graph of a successor function on \(V\).
\end{hbtheorem}

\begin{hbtheorem}{thm:p03-w03-p03-058-ordered}{cited=Schwentick1994}{Connectivity on ordered graphs}
The Connectivity query is not monadic ESO-definable on the class of finite
ordered graphs, that is, of finite structures \(\mathbf G=(V,E,\le)\) in
which \(E\) is a binary relation on \(V\) and \(\le\) is a linear order on
\(V\).
\end{hbtheorem}

These are Kolaitis's Theorems~5.10 and~5.11
\cite[p.~29]{Kolaitis2003ExpressivePowerNotes}. He states them without
proof and refers to de Rougemont \cite{deRougemont1987} and to Schwentick
\cite{Schwentick1994}; neither proof is reproduced here. The manuscript does
not say, for these two theorems, whether \emph{connected} means connected
through the symmetric closure of \(E\) or strongly connected when \(E\) is
not symmetric; \cref{thm:p03-w03-p03-058-connectivity} treats both
readings on graphs without built-in relations, and the two lemmas below
hold for both readings.

The hypotheses mark where the theorems stop.
\begin{enumerate}
\item \emph{The built-in relation strengthens the result.} A monadic
      ESO-sentence over \(\{E\}\) is also one over \(\{E,\mathrm{Suc}\}\)
      and over \(\{E,\le\}\), and it is evaluated on the reduct
      \((V,E)\). So each theorem implies the corresponding statement for
      the class of all finite graphs
      (\cref{thm:p03-w03-p03-058-reduct}). The converse is not automatic,
      because a defining sentence may use \(\mathrm{Suc}\) or \(\le\).
\item \emph{The chapter's strategy does not carry over.} The counting
      criterion \cref{thm:p03-w03-p03-058-hanf} is proved for coloured
      graphs with symmetric irreflexive edge relation only, and neighbourhood
      counting cannot replace it on ordered graphs. In the Gaifman graph of
      \((V,E,\le)\) any two distinct elements are adjacent, because one of
      them is \(\le\) the other \cite[Def.~3.22, p.~18]{Kolaitis2003ExpressivePowerNotes}.
      So every neighbourhood \(N(a,d)\) with \(d\ge2\) is the whole
      universe, as Kolaitis notes for linear orders
      \cite[Example~3.23, p.~18]{Kolaitis2003ExpressivePowerNotes}; the
      \(d\)-type of \(a\) is then the isomorphism type of the pointed
      structure \((\mathbf G,a)\), and two finite ordered graphs with an
      element of a common \(d\)-type are isomorphic. Kolaitis makes the same
      point for first-order definability on finite linear orders
      \cite[p.~20]{Kolaitis2003ExpressivePowerNotes}.
\item \emph{Monadic is essential.} Connectivity and strong connectivity are
      ESO-definable, by sentences whose quantified relation variables are
      all binary, on every class of finite structures with a binary
      relation \(E\) (\cref{thm:p03-w03-p03-058-binary-conn}). So both
      theorems, and \cref{thm:p03-w03-p03-058-connectivity}, do not extend
      to binary ESO. The displayed witness for connectivity quantifies one
      binary relation, while the witness for strong connectivity quantifies
      two. Problem~5.9 asks
      for a different query: one in \(\mathsf{NP}\) that no ESO-sentence
      with only binary relation variables defines.
\end{enumerate}

\begin{hblemma}{thm:p03-w03-p03-058-reduct}{proved=annex:p03-w03-p03-058-reduct}{Adding relations to the vocabulary}
Let \(\sigma\subseteq\sigma'\) be vocabularies, \(\mathcal C\) a class of
\(\sigma\)-structures, \(\mathcal C'\) a class of \(\sigma'\)-structures
whose \(\sigma\)-reducts all lie in \(\mathcal C\), \(Q\) a Boolean query on
\(\mathcal C\), and \(Q'(\mathbf G)=Q(\mathbf G{\restriction}\sigma)\) for
\(\mathbf G\in\mathcal C'\). If \(Q\) is (monadic) ESO-definable on
\(\mathcal C\), then \(Q'\) is (monadic) ESO-definable on \(\mathcal C'\) by
the same sentence. In particular,
\cref{thm:p03-w03-p03-058-successor} and
\cref{thm:p03-w03-p03-058-ordered} each imply that the Connectivity query,
in the reading used there, is not monadic ESO-definable on the class of all
finite graphs.
\end{hblemma}

\begin{hblemma}{thm:p03-w03-p03-058-binary-conn}{proved=annex:p03-w03-p03-058-binary-conn}{Connectivity in binary ESO}
Let \(\sigma\) be a vocabulary containing a binary relation symbol \(E\),
and let \(\lambda(L)\) be the first-order formula saying that \(L\) is a
strict linear order (irreflexive, transitive and total). Put
\[
 \gamma=(\exists L)\Bigl(\lambda(L)\wedge(\forall y)\bigl((\exists w)L(w,y)
 \to(\exists z)(L(z,y)\wedge(E(z,y)\vee E(y,z)))\bigr)\Bigr)
\]
and
\[
 \begin{aligned}
 \gamma_s=(\exists L_1)(\exists L_2)\Bigl(&\lambda(L_1)\wedge\lambda(L_2)
 \wedge(\forall x)\bigl(\neg(\exists w)L_1(w,x)\leftrightarrow
 \neg(\exists w)L_2(w,x)\bigr)\\
 &\wedge(\forall y)\bigl((\exists w)L_1(w,y)\to(\exists z)(L_1(z,y)\wedge
 E(z,y))\bigr)\\
 &\wedge(\forall y)\bigl((\exists w)L_2(w,y)\to(\exists z)(L_2(z,y)\wedge
 E(y,z))\bigr)\Bigr).
 \end{aligned}
\]
On every finite \(\sigma\)-structure, \(\gamma\) holds exactly when
\((V,E)\) is connected through the symmetric closure of \(E\), and
\(\gamma_s\) holds exactly when \((V,E)\) is strongly connected. Both
sentences quantify only binary relation variables.
\end{hblemma}

\emph{Relationships.} With \(k=0\) relations,
\cref{thm:p03-w03-p03-058-eso-method} reduces to the first-order game
method, \cref{thm:p03-w03-p03-055-rank-definability}, and the cycles
\(C_{2m}\) and \(C_m\sqcup C_m\) used there for first-order connectivity
(\cref{thm:p03-w03-p03-038-connectivity}) reappear here inside a game in which
Spoiler may first colour the graph. Fagin's theorem (\cref{thm:p03-fagin})
identifies ESO-definability on finite structures with membership in
\(\mathsf{NP}\), which is what lets \cref{thm:p03-w03-p03-058-np-conp}
speak about complexity classes. The independence-friendly evaluation game
of \cref{def:p03-if-logic} is a different game
(\cref{def:p03-w03-p03-058-not-if}).
